% !TeX root = ../papers/nc-p3.tex
\section{極をもつ傾きでの八重被覆と六次の積表}
\label{sec:rank-fifteen-pole-fibers}
\begin{theorem}\label{thm:rank-fifteen-pole-fibers}
剰余標数3の標数零の通常な階数15窓で $\nu(x)=0$, $\bar x\ne1$ とする。
$0<3\nu(y)\le2\nu(3)$ なら、各始点で四つの中心四次式上自由階数256の
標準Koszul左右AS正則窓である。

\end{theorem}
\begin{proof}
定理\ref{thm:rank-fifteen-pole-slope-annulus}の表示で
$\nu(r)=\nu(y)/4$、$x=(4+r^2)(AB)^2$ である。
$s$ から剰余座標 $X=\bar x$ への写像は
\[
  X=N(s)/D(s),\qquad N=16s^2(1-s^2)^2,\qquad D=(1+s^2)^4.
\]
互いに素な二式による次数8の写像で、$s\mapsto-s,1/s,(1+s)/(1-s)$ が生成する
八つの相異なる分数線形変換はこれを保つ。実際、第三の変換の二乗は $-1/s$ であり、
$(A,B)$ の符号変更と交換がこれらの表示を与える。従って一般のファイバーは一軌道で、
$X\ne0,1,\infty$ では $A,B,A+B,A-B$ がいずれも非零で八表示は相異なる。
各表示は符号を整えれば所定の窓へ持ち上がる。実際 $T=AB$ として
定理\ref{thm:rank-fifteen-pole-slope-annulus}の $T$ の平方根の符号を選び、
$T(1+s^2)^2-2s(1-s^2)=0$ の指定した剰余根を持ち上げればよい。

\path{computations/verify_rank_fifteen_pole_fibers.py} は、同じ十七桁の
$\mathbb F_3(s)$ 上の計算で、単元として用いる主元の定数係数の分子と分母を記録する。
平方自由な積の次数は121である。$X=0,1,\infty$ 上の因子を除き、
多項式と上の三変換による引き戻しの最大公約元を安定するまで取る。
モニック化した結果は次数48であり、
\[
  D^6 g_1(N/D)g_2(N/D),\qquad
  g_1(X)=X^3+2X^2+2X+2,\quad g_2(X)=X^3+X^2+2
\]
のモニック化に等しい。従ってこれらの根を除く各ファイバーに良い表示がある。
二つの三次式の根ではそれぞれ $s^3+s^2+s+2$,
$s^6+2s^4+2s^3+s^2+2$ の根を選ぶ。対応する二つの剰余体で同時に
単元となる主元を選んでも、除算前の消滅は $\mathbb F_3(s)$ 上で成り立つ。
どちらの表示でも中心商の長さ6から13での最大損失は $(2,1,1,1,1,0,1,0)$ である。

一方、基準表示で $\mathbb Q(s)$ 上の十二桁を使い、
未切断経路、四中心、五次の積表に続いて六次の全十六行列を消去する。
全ての $r$ 除算で捨てる係数はこの有理関数体で厳密に零である。
主元の先頭の分子・分母を3で還元した平方自由な積は正確に $s^9-s$ であり、
$ND(N-D)$ を割る。従って $X\ne0,1,\infty$ のどの表示でも
六次までの同じ消去を混標数の形式級数環で行える。
六次の最大損失は二桁で、一桁を残して文字別 $(12,12,10,10)$ の基底を得る。
中心の自由座標とそれを指定する四次語基底は、標数3の基準表示と追加表示の双方で一致するため、六次までの自由な語の格子の左右の積を
標数3に還元し、同じ商を得る。十七桁の剰余計算では六次の後に六桁が残り、
その後の最大損失の和は五である。従って $3\in r^6\mathcal O$ なら
六次から再開して十三次まで進める。十三次の消滅まで全次元上界が $(1+t+t^2+t^3)^4$ の係数に一致するので、
第\ref{sec:rank-fifteen-overlap-annulus}節の自由性と左右の解消の議論が適用できる。
\end{proof}
